% Chapter 1, Topic from _Linear Algebra_ Jim Hefferon
%  http://joshua.smcvt.edu/linearalgebra
%  2001-Jun-09
\topic{Analyzing Networks}
\index{networks|(}
This diagram shows an electrical network.
The battery is on the left, drawn as stacked line segments. 
The wires are lines, shown straight and with sharp right angles
for neatness.
The zigzag's are \definend{resistors}.\index{resistor}
\begin{center}
  \includegraphics{gr/mp/ch1.42}
\end{center}
Network designers need to know how 
much electricity flows 
through each component.
We will use linear systems to analyze simple 
electrical networks. 
We need two facts about electricity 
and two facts about networks.

The first electrical fact is that a battery is like a pump,  
providing a force impelling the electricity to flow.  
We say that the battery provides a 
\definend{potential}.\index{potential, electric}
That force creates a
flow through the circuit, called a \definend{current}.

The second fact is that in resistors
the amount of flow is proportional to the force provided by the battery.
For each such component there is a number,  
its \definend{resistance},\index{resistance}
such that the potential is equal to the flow times the resistance.
Potential is measured in \definend{volts},
the flow is in \definend{amperes},
and resistance to that flow is in \definend{ohms};
these units are defined so that
$\mbox{volts}=\mbox{amperes}\cdot\mbox{ohms}$.

An example is that if a resistor measures $2$~ohms 
then wiring it to a $12$~volt battery 
results in a flow of $6$~amperes.
Conversely, if electrical current of
$2$~amperes flows through that resistor then there must be 
a $4$~volt potential difference between its ends. 
This is the \definend{voltage drop}\index{voltage drop} across the 
resistor.
One way to think of the electrical circuits that we consider here
is that the battery provides a voltage rise while the other components 
are voltage drops.

The facts that we need about networks are  
\textit{Kirchoff's Current Law}, that for any point in a network the flow in
     equals the flow out and
\textit{Kirchoff's Voltage Law}, that around any circuit the total drop equals 
     the total rise.\index{Kirchhoff's Laws}\index{networks!Kirchhoff's Laws}
% In the above network there is only one voltage rise, at the battery, but
% some networks have more than one.

We start with the simple network below.
It has a battery and three resistors.
When components are wired one after another, as here,
they are in \definend{series}.\index{circuits!series}
\begin{center}
  \includegraphics{gr/mp/ch1.35}
\end{center}
By Kirchoff's Voltage Law, because the voltage rise is
$20$~volts, the total voltage drop must also be $20$~volts.
Since the resistance from start to finish is
$10$~ohms (the resistance of the wire connecting the components is negligible),
the current is $(20/10)=2$~amperes. 
Now, by Kirchhoff's Current Law, there are $2$~amperes through
each resistor.
Therefore the voltage drops are: 
$4$~volts across the $2$~ohm resistor,
$10$~volts across the $5$~ohm resistor, 
and $6$~volts across the $3$~ohm resistor.

The prior network is simple enough that we didn't use a linear system, but
the next one is more complicated.
Here the resistors are in \definend{parallel}.\index{circuits!parallel}
\begin{center}
  \includegraphics{gr/mp/ch1.36}
\end{center}
We begin by labeling the branches as below.
Let the current through the left branch of the parallel portion
be $i_1$ and that through
the right branch be $i_2$, and also
let the current through the battery be $i_0$.
We don't need to know 
the actual direction of flow\Dash if current flows in the direction 
opposite to our arrow 
then we will get a negative number in the solution.
\begin{center}
  \includegraphics{gr/mp/ch1.37}
\end{center}
The Current Law, applied to the split point in the upper right, 
gives that $i_0=i_1+i_2$. 
Applied to the split point in the lower right it gives $i_1+i_2=i_0$.  
In the circuit that loops out of the top of the battery, 
down the left branch of the
parallel portion, and back into the bottom of the battery, the voltage
rise is $20$ while the voltage drop is $i_1\cdot 12$, so
the Voltage Law gives that $12i_1=20$.
Similarly, the circuit from the battery to the right branch and back to the
battery gives that $8i_2=20$.
And, in the circuit that simply loops around in 
the left and right branches of the parallel portion 
(we arbitrarily take the direction of clockwise), 
there is a voltage rise of $0$ and a voltage drop of $8i_2-12i_1$
so $8i_2-12i_1=0$.
\begin{displaymath}
  \begin{linsys}{3}
   i_0&- &i_1    &-  &i_2   &=  &0 \\
  -i_0&+ &i_1    &+  &i_2   &=  &0  \\
      &  &12i_1  &   &      &=  &20  \\
      &  &       &   &8i_2  &=  &20  \\
      &  &-12i_1 &+  &8i_2  &=  &0  
  \end{linsys}
\end{displaymath}
The solution
is $i_0=25/6$, $i_1=5/3$, and $i_2=5/2$, all in amperes.
(Incidentally, the redundant equations illustrate 
that they can arise in practice.) 

Kirchhoff's laws can 
establish the electrical properties of very complex networks.
The next diagram shows
five resistors, whose values are in ohms, wired in 
\definend{series-parallel}.\index{circuits!series-parallel} 
\begin{center}
  \includegraphics{gr/mp/ch1.38}
\end{center}
This is a \definend{Wheatstone bridge}\index{Wheatstone bridge} 
(see \nearbyexercise{exer:WheatstoneBr}).
To analyze it, we can place the arrows in this way.
\begin{center}
  \includegraphics{gr/mp/ch1.39}
\end{center}
Kirchhoff's Current Law, applied to the
top node, the left node, the right node, and the bottom node gives
these.
\begin{align*}
   i_0     &=  i_1+i_2  \\
   i_1     &=  i_3+i_5  \\
   i_2+i_5 &=  i_4      \\
   i_3+i_4 &=  i_0
\end{align*} 
Kirchhoff's Voltage Law,
applied to the inside loop (the $i_0$ to~$i_1$ to~$i_3$ to~$i_0$ loop), 
the outside loop, 
and the upper loop not involving the battery, gives these.  
\begin{align*}
      5i_1+10i_3  &= 10   \\
      2i_2+4i_4   &= 10   \\
      5i_1+50i_5-2i_2  &= 0     
\end{align*} 
Those suffice to determine the solution 
$i_0=7/3$, $i_1=2/3$, $i_2=5/3$, 
$i_3=2/3$, $i_4=5/3$, and $i_5=0$. 

We can understand many kinds of networks in this way.
For instance, the exercises analyze some networks of streets.
 
\begin{exercises}
  % \item[] \textit{Many of the systems for these problems
  %     are mostly easily solved on a computer.}
  \item 
    Calculate the amperages in each part of each network.
    \begin{exparts}
      \partsitem This is a simple network.
          \begin{center}
             \includegraphics{gr/mp/ch1.40}
        \end{center}
      \partsitem Compare this one with the parallel case discussed above.
          \begin{center}
             \includegraphics{gr/mp/ch1.41}
        \end{center}
      \partsitem We used this network to introduce this topic.
          \begin{center}
           \includegraphics{gr/mp/ch1.42}
        \end{center}
    \end{exparts}
    \begin{answer}
      \begin{exparts}
        \partsitem The total resistance is $7$~ohms.
          With a $9$~volt potential, the flow will be $9/7$~amperes.
          The voltage drops will then be:~$27/7$~volts
          across the $3$~ohm resistor, and $18/7$~volts across each of
          the two $2$~ohm resistors.        
        \partsitem The $2$~ohm
          resistor on the left has a voltage drop of $9$~volts
          (and hence the flow through it is $9/2$~amperes), and the 
          remaining portion on the right also has a voltage drop of 
          $9$~volts, and so we can analyze it as in the prior item.
          We can also use linear systems.
          \begin{center}
            \includegraphics{gr/mp/ch1.48}
          \end{center}
          The Current Law applied to junction in the the upper middle gives 
          $i_0=i_1+i_2$.
          Applied to the lower middle junction it gives $i_1=i_2+i_3$.
          The Voltage Law applied to the leftmost circuit
          is that the $9$~volts up equals the voltage drop across the
          $2$~ohm resistor, $9=2i_1$.
          The rightmost circuit has a total resistance of $3+2+2=7$~ohm,
          so the Voltage Law gives $9=7i_2$. 
          Here is the linear system.
          \begin{equation*}
            \begin{linsys}{4}
              i_0  &- &i_1  &- &i_2  &  &    &= &0  \\
                   &  &i_1  &+ &i_2  &- &i_3 &= &0  \\
                   &  &2i_1 &  &     &  &    &= &9  \\  
                   &  &     &  &7i_2 &  &    &= &9       
            \end{linsys}
          \end{equation*}
          Gauss's Method
          \begin{multline*}
            \grstep{-2\rho_{2}+\rho_{3}}
            \begin{amat}{4}
              1  &-1  &-1  &0  &0  \\ 
              0  &1  &1  &-1  &0  \\ 
              0  &0  &-2  &2  &9  \\ 
              0  &0  &7  &0  &9  \\ 
            \end{amat}                          
            \grstep{(7/2)\rho_{3}+\rho_{4}}
            % \begin{amat}{4}
            %     1  &-1  &-1  &0  &0  \\ 
            %     0  &1  &1  &-1  &0  \\ 
            %     0  &0  &-2  &2  &9  \\ 
            %     0  &0  &0  &7  &81/2  \\ 
            %   \end{amat}                              \\
            \grstep[(1/7)\rho_{4}]{(-1/2)\rho_{3}}         \\
            % \begin{amat}{4}
            %   1  &-1  &-1  &0  &0  \\ 
            %   0  &1  &1  &-1  &0  \\ 
            %   0  &0  &1  &-1  &-9/2  \\ 
            %   0  &0  &0  &1  &81/14  \\ 
            % \end{amat}
            \grstep[\rho_{4}+\rho_{3}]{\rho_{4}+\rho_{2}}
            % \begin{amat}{4}
            %   1  &-1  &-1  &0  &0  \\ 
            %   0  &1  &1  &0  &81/14  \\ 
            %   0  &0  &1  &0  &9/7  \\ 
            %   0  &0  &0  &1  &81/14  \\ 
            % \end{amat}                    \\
            \grstep[-\rho_{3}+\rho_{2}]{\rho_{3}+\rho_{1}}   
            % \begin{amat}{4}
            %   1  &-1  &0  &0  &9/7  \\ 
            %   0  &1  &0  &0  &9/2  \\ 
            %   0  &0  &1  &0  &9/7  \\ 
            %   0  &0  &0  &1  &81/14  \\ 
            % \end{amat}
            \grstep{\rho_{2}+\rho_{1}}
            \begin{amat}{4}
              1  &0  &0  &0  &81/14  \\ 
              0  &1  &0  &0  &9/2  \\ 
              0  &0  &1  &0  &9/7  \\ 
              0  &0  &0  &1  &81/14  \\ 
            \end{amat}
          \end{multline*}
          gives the unique solution $i_0=81/14$, $i_1=9/2$, $i_2=9/7$,
          and $i_3=81/14$.
        \partsitem
          Using these variables,
          \begin{center}
            \includegraphics{gr/mp/ch1.49}
          \end{center}
          below is one linear system sufficient to yield the solution.
          The first four equations are from the Current Law.
          The last three are from applying the Voltage Law to 
          the circuit involving
          $i_0$-$i_1$-$i_6$, 
          the circuit involving $i_0$-$i_2$-$i_4$-$i_5$-$i_6$, 
          and the circuit with $i_0$-$i_2$-$i_3$-$i_5$-$i_6$.
          \begin{equation*}
            \begin{linsys}{7}
              i_0  &- &i_1  &- &i_2  &  &    &  &    & &    & &    &= &0  \\
                   &  &     &  &i_2  &- &i_3 &- &i_4 & &    & &    &= &0  \\
                   &  &     &  &     &  &i_3 &+ &i_4 &-&i_5 & &    &= &0  \\  
                   &  &i_1  &  &     &  &    &  &    &+&i_5 &-&i_6 &= &0  \\  
                   &  &3i_1 &  &     &  &    &  &    & &    & &    &= &9  \\  
                   &  &     &  &3i_2 &  &    &+ &2i_4&+&2i_5& &    &= &9  \\  
                   &  &     &  &3i_2 &+ &9i_3&  &    &+&2i_5& &    &= &9  
            \end{linsys}
          \end{equation*}
          Gauss's Method gives
          $i_0=318/73\approx 4.36$, 
          $i_1=3$, 
          $i_2=99/73\approx 1.36$,
          $i_3=18/73\approx 0.25$, 
          $i_4=81/73\approx 1.11$, 
          $i_5=99/73\approx 1.36$, 
          and $i_6=318/73\approx 4.36$.
      \end{exparts}
    \end{answer}
  \item 
    In the first network that we analyzed, with the three resistors  
    in series, we just added to get
    that they acted together like a single resistor of $10$~ohms.
    We can do a similar thing for parallel circuits. 
    In the second circuit analyzed,
    \begin{center}
      \includegraphics{gr/mp/ch1.36}
    \end{center}
    we found that
    the current through the battery is $25/6$~amperes.
    Thus, the parallel portion is 
    \definend{equivalent}\index{resistance!equivalent}
    to a single resistor of 
    $20/(25/6)=4.8$~ohms. 
    \begin{exparts}
      \partsitem What is the equivalent resistance if we change
         the $12$~ohm resistor to $5$~ohms?
      \partsitem What is the equivalent resistance if the two are each
         $8$~ohms?
      \partsitem Find the formula for the equivalent resistance if
         the two resistors in parallel are $r_1$~ohms and $r_2$~ohms.
    \end{exparts}
    \begin{answer}
      \begin{exparts}
        \partsitem 
          Using the variables from the earlier analysis,
          \begin{displaymath}
            \begin{linsys}{3}
              i_0&- &i_1    &-  &i_2   &=  &0 \\
             -i_0&+ &i_1    &+  &i_2   &=  &0  \\
                 &  &5i_1   &   &      &=  &20  \\
                 &  &       &   &8i_2  &=  &20  \\
                 &  &-5i_1  &+  &8i_2  &=  &0  
            \end{linsys}
          \end{displaymath}
          Gauss's Method gives that the current flowing in each branch 
          is $i_0= 13/2=6.5$,
          $i_1= 4$, 
          and $i_2= 5/2=2.5$, 
          all in amperes.
          Thus the parallel portion is acting as a single resistor
          of $20/(13/2)\approx 3.08$~ohms.
        \partsitem 
          A similar analysis 
          \begin{displaymath}
            \begin{linsys}{3}
              i_0&- &i_1    &-  &i_2   &=  &0 \\
             -i_0&+ &i_1    &+  &i_2   &=  &0  \\
                 &  &8i_1   &   &      &=  &20  \\
                 &  &       &   &8i_2  &=  &20  \\
                 &  &-8i_1  &+  &8i_2  &=  &0  
            \end{linsys}
          \end{displaymath}
          gives that
          $i_o=5$,  $i_1=5/2=2.5$, and $i_2=5/2=2.5$~amperes.
          The equivalent resistance is $20/5=4$~ohms.
        \partsitem 
          Another analysis like the prior ones 
          \begin{displaymath}
            \begin{linsys}{3}
              i_0&- &i_1      &-  &i_2     &=  &0 \\
             -i_0&+ &i_1      &+  &i_2     &=  &0  \\
                 &  &r_1i_1   &   &        &=  &20  \\
                 &  &         &   &r_2i_2  &=  &20  \\
                 &  &-r_1i_1  &+  &r_2i_2  &=  &0  
            \end{linsys}
          \end{displaymath}
          gives $i_2=20/r_2$, $i_1=20/r_1$, 
          and $i_0=20(r_1+r_2)/(r_1r_2)$, all in amperes.
          The parallel portion is acting as a single resistor of
          size $20/i_0=r_1r_2/(r_1+r_2)$~ohms.
          (This equation is often stated as:~the equivalent 
          resistance~$r$ satisfies $1/r=(1/r_1)+(1/r_2)$.)
      \end{exparts}
    \end{answer}
\item \label{exer:WheatstoneBr} 
   A Wheatstone bridge\index{Wheatstone bridge}
   is an instrument that measures resistance. 
   \begin{center}
     \includegraphics{gr/mp/ch1.47}
   \end{center}
   Show that in this circuit if the current
   flowing through $r_g$ is zero then $r_4=r_2r_3/r_1$.
   (To operate the device, put the unknown
   resistance at $r_4$.
   At $r_5$ is a meter that shows the current.
   We vary the three resistances $r_1$, $r_2$, and $r_3$, which
   they each have a calibrated knob, until the 
   current in the middle reads $0$.
   Then the equation gives the value of $r_4$.)
   \begin{answer}
     Here is the graphic from the text body that establishes the sign convention
     for the current flows.
     \begin{center}
       \includegraphics{gr/mp/ch1.39}
     \end{center}
     Apply Kirchoff's Current Law to the node where 
     $i_1$, $i_3$, and~$i_5$ come together, and also 
     to the node where $i_2$, $i_4$, and~$i_5$ come
     together.
     \begin{equation*}
       \begin{linsys}{3}
         i_1 &- &i_3 &- &i_5 &= &0  \\
         i_2 &- &i_4 &+ &i_5 &= &0  
       \end{linsys}
     \end{equation*}
     Kirchoff's Voltage law, applied to the loop in the top right, 
     and to the loop in the bottom right, gives these. 
     \begin{equation*}
       \begin{linsys}{3}
         i_2r_2 &- &i_5r_5 &- &i_1r_1 &= &0  \\
         i_4r_4 &- &i_3r_3 &+ &i_5r_5 &= &0  
       \end{linsys}
     \end{equation*}
     Assume that $i_5$ is zero. 
     Then $i_1=i_3$, and $i_2=i_4$,
     and $i_1r_1=i_2r_2$, and $i_3r_3=i_4r_4$.
     Rearrange the last equality
     \begin{equation*}
       r_4=\frac{i_3r_3}{i_4}\cdot\frac{i_2r_2}{i_1r_1}
     \end{equation*}
     and cancel the $i$'s to get the desired conclusion.
   \end{answer}
   %  \item Prove Th\`evenin's Theorem.
% \item[]\textit{There are networks other than electrical ones, and 
%               we can ask how well Kirchhoff's laws apply to them.
%               The remaining questions consider an extension to 
%               networks of streets.}
\item 
    Consider this traffic circle.
    \begin{center}
      \includegraphics{gr/mp/ch1.43}
     \end{center}
     This is the traffic volume, in units of cars per ten minutes.    
     \begin{center}
       \begin{tabular}{r|ccc}
        \multicolumn{1}{c}{\ } % get rid of vert bar
                            &\textit{North}  
                            &\textit{Pier}  
                            &\textit{Main}  \\
         \cline{2-4}
         \begin{tabular}{@{}r@{}} \textit{into} \\ \textit{out of} \end{tabular}
         &\begin{tabular}{@{}r@{}} $100$ \\ $75$ \end{tabular}
         &\begin{tabular}{@{}r@{}} $150$ \\ $150$ \end{tabular}
         &\begin{tabular}{@{}r@{}} $25$ \\  $50$ \end{tabular}
         % rows:
         % \textit{into}      &$100$           &$150$          &$25$      \\
         % \textit{out of}    &$75$            &$150$          &$50$       
       \end{tabular}
    \end{center}
    We can set up equations to model how the traffic flows around the circle.   
    \begin{exparts}
      \partsitem 
         Adapt Kirchhoff's Current Law to this circumstance.
         Is it a reasonable modeling assumption?
      \partsitem 
         Let 
         $i_1$ be the rate on the circle section going 
         counterclockwise from North Avenue to Main,
         in number of cars per ten minutes.
         Let $i_2$ be the counterclockwise rate on the circle section between
         Main and Pier,
         and let $i_3$ be the rate between Pier and North.
         Using the adapted law,
         for each of the three circle sections state an equation
         describing the traffic flow at that node.
      \partsitem 
         Solve that system.
      \partsitem
         Interpret your solution.
      \partsitem
         Restate the Voltage Law for this circumstance.
         How reasonable is it?
    \end{exparts}
    \begin{answer}
      \begin{exparts}
        \partsitem
           The natural adaptation is:~for any circle section, 
           the flow in equals the flow out.
           It seems reasonable in this case, unless cars are stuck 
           for a long time.
        \partsitem
           If we know $i_1$ then we can determine $i_2$ because between
           the two is Main, and $50$ cars leave via Main while $25$~cars enter. 
           Thus $i_1-25=i_2$.
           Similarly Pier's in/out balance means that $i_2=i_3$ and
           North gives $i_3+25=i_1$. 
           We have this system.
           \begin{equation*}
             \begin{linsys}{3}
               i_1  &-  &i_2  &   &     &=  &25  \\
                    &   &i_2  &-  &i_3  &=  &0   \\
              -i_1  &   &     &+  &i_3  &=  &-25
             \end{linsys}
           \end{equation*}
        \partsitem 
           Gauss's Method
           \begin{align*}
             \begin{amat}{3}
               1  &-1  &0  &25  \\ 
               0  &1  &-1  &0  \\ 
               -1  &0  &1  &-25  \\ 
             \end{amat}
             &\grstep{\rho_{1}+\rho_{3}}
             \begin{amat}{3}
               1  &-1  &0  &25  \\ 
               0  &1  &-1  &0  \\ 
               0  &-1  &1  &0  \\ 
             \end{amat}                       \\
             &\grstep{\rho_{2}+\rho_{3}}
             \begin{amat}{3}
               1  &-1  &0  &25  \\ 
               0  &1  &-1  &0  \\ 
               0  &0  &0  &0  \\ 
             \end{amat}
           \end{align*}
           gives that there are infinitely many solutions.
           With $i_3$ as the parameter, this is the solution set.
           \begin{equation*}
             \set{\colvec[c]{25+i_3  \\ i_3  \\i_3} \suchthat i_3\in\Re}
           \end{equation*}
        \partsitem
           If we picture an initially-empty circle with the given input/output
           behavior, we can superimpose $i_3$-many circling cars 
           to get a new solution.
        \partsitem
           A suitable restatement is that the number of cars entering the 
           circle must equal the number of cars leaving.
           This is reasonable over a long time period.
           Over one minute it might not be true but over a year, we expect
           it is an excellent approximation.
      \end{exparts}
    \end{answer}
  \item 
     This is a network of streets. 
          \begin{center}
            \includegraphics{gr/mp/ch1.44}
             \end{center}
         We can observe the hourly flow of cars into this network's
         entrances, and out of its exits.
          \begin{center}
            \begin{tabular}{r|ccccc}
               \multicolumn{1}{c}{}  
                 &\textit{east Winooski}  
                 &\textit{west Winooski}  
                 &\textit{Willow}  
                 &\textit{Jay} 
                 &\textit{Shelburne} \\ 
                 \cline{2-6}
              \begin{tabular}{@{}r@{}} \textit{into} \\ \textit{out of} \end{tabular} 
              &\begin{tabular}{@{}r@{}} $80$ \\ $30$ \end{tabular} 
              &\begin{tabular}{@{}r@{}} $50$ \\ $5$ \end{tabular} 
              &\begin{tabular}{@{}r@{}} $65$ \\ $70$ \end{tabular} 
              &\begin{tabular}{@{}r@{}} -- \\ $55$ \end{tabular} 
              &\begin{tabular}{@{}r@{}} $40$ \\ $75$ \end{tabular} 
              % rows:
              % \textit{into}      &80    &50    &65     &--    &40      \\
              % \textit{out of}    &30    &5     &70     &55    &75      
            \end{tabular}
          \end{center}
          To reach Jay a
          car must enter the network via some other road first, which is why
          there is no `into Jay' entry in the table.
          Also, over a long time 
          the total in must approximately equal the total 
          out, which is why both rows add to $235$~cars.
          Once inside the network, the traffic may flow in different
          ways, perhaps filling Willow and leaving Jay 
          mostly empty, etc.
          Kirchhoff's Laws limit that freedom.
    \begin{exparts}
       \partsitem Determine the restrictions on the flow inside this network 
          of streets by setting
          up a variable for each block, establishing the equations, 
          and solving them.
          Notice that some streets are one-way only.
          \textit{Hint:}~On the map above, use the data from the table to
          label the flows on some of the blocks.
          Name each remaining block with a variable, and for each intersection
          between two streets get a relationship between some of these variables.
       \partsitem Suppose that someone proposes construction for
         Winooski Avenue East between Willow and Jay, 
         and traffic on that block will be reduced.
         What is the least amount of traffic flow that can we can
         allow on that block without disrupting the  
         hourly flow into and out of the network?
    \end{exparts}
    \begin{answer}
      \begin{exparts}
        \partsitem Here the blocks with flows known from the table are 
          labeled, and for other blocks we declare a variable.
          \begin{center}
            \includegraphics{gr/mp/ch1.51}          
          \end{center}
          Consider the intersection Willow and Shelburne, in the upper left. 
          Kirchhoff's principle that the flow in must equal the flow out gives
          $i_1+65=i_2+75$, and so $i_1-i_2=10$.
          Similarly at Shelburne and Jay we get $40+i_3=i_1+55$. 
          Doing the intersections from right to left and top to bottom
          gives this system.
          \begin{equation*}
            \begin{linsys}{7}
              i_1 &- &i_2 &  &    &  &    &  &    &  &    &  &    &= &10  \\
             -i_1 &  &    &+ &i_3 &  &    &  &    &  &    &  &    &= &15   \\
                  &  &i_2 &  &    &+ &i_4 &- &i_5 &  &    &  &    &= &5   \\
                  &  &    &  &-i_3&- &i_4 &  &    &+ &i_6 &  &    &= &-50 \\
                  &  &    &  &    &  &    &  &i_5 &  &    &- &i_7 &= &-10 \\
                  &  &    &  &    &  &    &  &    &  &-i_6&+ &i_7 &= &30 
            \end{linsys}
          \end{equation*}
          The row operation $\rho_1+\rho_2$ followed by $\rho_2+\rho_3$
          then $\rho_3+\rho_4$ and $\rho_4+\rho_5$ and finally $\rho_5+\rho_6$
          result in this system.
          % runfile("gauss_method.sage")
          % M = matrix(QQ, [[1,-1,0,0,0,0,0], [-1,0,1,0,0,0,0], [0,1,0,1,-1,0,0], [0,0,-1,-1,0,1,0], [0,0,0,0,1,0,-1], [0,0,0,0,0,-1,1]])
          % v = vector(QQ, [10,15,5,-50,-10,30])
          % M_prime = M.augment(v, subdivide=True)
          % gauss_method(M_prime)
          \begin{equation*}
            \begin{linsys}{7}
              i_1 &- &i_2 &  &    &  &    &  &    &  &    &  &    &= &10  \\
                  &  &-i_2&+ &i_3 &  &    &  &    &  &    &  &    &= &25   \\
                  &  &    &  &i_3 &+ &i_4 &- &i_5 &  &    &  &    &= &30  \\
                  &  &    &  &    &  &    &  &-i_5&+ &i_6 &  &    &= &-20  \\
                  &  &    &  &    &  &    &  &    &  &i_6 &- &i_7 &= &-30 \\
                  &  &    &  &    &  &    &  &    &  &    &  &0   &= &0   
            \end{linsys}
          \end{equation*}
          Take the free variables $i_4$ and~$i_7$ as 
          parameters.
          \begin{equation*}
          \begin{split}
            i_6  &=  i_7-30  \\
            i_5  &=  i_6+20=(i_7-30)+20=i_7-10 \\
            i_3  &=  -i_4+i_5+30=-i_4+(i_7-10)+30=-i_4+i_7+20 \\
            i_2  &=  i_3-25=(-i_4+i_7+20)-25=-i_4+i_7-5 \\
            i_1  &=  i_2+10=(-i_4+i_7-5)+10=-i_4+i_7+5
          \end{split}
          \end{equation*}
          Obviously $i_4$ and $i_7$ must be positive, and in fact
          the first equation shows that $i_7$ must be at least $30$.
        \partsitem We cannot take $i_7$ to be zero or else $i_6$ will
          be negative (this would mean cars going the wrong way on the
          one-way street Jay).
          We can, however, take $i_7$ to be as small as $30$, and then 
          there are many suitable $i_4$'s.
          For instance, the solution
          \begin{equation*}
            (i_1,i_2,i_3,i_4,i_5,i_6,i_7)
            =
            (35,25,50,0,20,0,30)
          \end{equation*}
          results from choosing $i_4=0$.
      \end{exparts}
    \end{answer}
\end{exercises}
\index{networks|)}
\endinput
